\chapter[Sharding y Distribución]{\emj{🌐}~Sharding y Distribución}

\section[Cuándo escalar horizontalmente]{\emj{📈}~Cuándo escalar horizontalmente}

Escalar verticalmente (más CPU, RAM, SSD) es siempre la primera respuesta: más simple, sin cambios de arquitectura.
El sharding horizontal se justifica cuando una sola máquina no puede manejar el volumen de escrituras o el tamaño de datos supera lo manejable en un nodo.

\begin{teoriabox}{Límites que fuerzan la distribución}
\begin{itemize}
  \item \textbf{Write throughput:} un primary puede saturar su I/O de WAL. Sharding divide las escrituras entre múltiples primaries.
  \item \textbf{Dataset size:} terabytes que no caben en RAM de un nodo degradan el performance a I/O puro.
  \item \textbf{Geografía:} usuarios en distintas regiones necesitan latencia baja → datos cerca de ellos.
  \item \textbf{Compliance:} datos de ciertos países deben residir en esa jurisdicción.
\end{itemize}
\textbf{Costo del sharding:} queries cross-shard se vuelven joins distribuidos (caros), transacciones distribuidas (2PC = complejo), y el operacional se multiplica.
\end{teoriabox}

\section[Partitioning nativo de PostgreSQL]{\emj{🗂️}~Partitioning nativo de PostgreSQL}

El partitioning es el primer paso antes de sharding real. Divide una tabla lógica en múltiples tablas físicas (particiones) en el \emph{mismo} nodo.

El partitioning declarativo de PostgreSQL (disponible desde PG 10, mejorado significativamente en PG 11--14) divide una tabla lógica en múltiples tablas físicas llamadas particiones.
Para el código de la aplicación, la tabla padre se ve como una tabla normal: los INSERTs, UPDATEs, SELECTs se escriben contra la tabla padre y PG los enruta a la partición correcta automáticamente.
El beneficio principal es \textbf{partition pruning}: el planner elimina de la ejecución las particiones que no contienen datos relevantes para la query, reduciendo dramáticamente el I/O en tablas con datos temporales o por categoría.

\begin{lstlisting}[caption={Table Partitioning por rango (tiempo)}]
-- Tabla madre particionada:
CREATE TABLE events (
  id          BIGSERIAL,
  user_id     BIGINT NOT NULL,
  event_type  TEXT NOT NULL,
  created_at  TIMESTAMPTZ NOT NULL DEFAULT NOW(),
  payload     JSONB
) PARTITION BY RANGE (created_at);

-- Particiones por mes:
CREATE TABLE events_2025_01 PARTITION OF events
  FOR VALUES FROM ('2025-01-01') TO ('2025-02-01');

CREATE TABLE events_2025_02 PARTITION OF events
  FOR VALUES FROM ('2025-02-01') TO ('2025-03-01');

-- Índices se crean por partición:
CREATE INDEX ON events_2025_01 (user_id);
CREATE INDEX ON events_2025_02 (user_id);

-- Partición default para fuera de rango:
CREATE TABLE events_default PARTITION OF events DEFAULT;

-- PG rutea automáticamente INSERT/UPDATE/DELETE a la partición correcta
INSERT INTO events (user_id, event_type, created_at)
VALUES (42, 'click', '2025-01-15');
-- Automáticamente va a events_2025_01
\end{lstlisting}

El partition pruning solo funciona si la condición WHERE usa la \emph{misma expresión} que define la partición.
Si la columna de partición es \texttt{created\_at} y filtras por \texttt{DATE(created\_at)}, el planner no puede determinar qué particiones eliminar y escanea todas.
Siempre verificar con EXPLAIN que solo aparecen las particiones esperadas.

\begin{lstlisting}[caption={Partition pruning — verificar que funciona}]
-- Con EXPLAIN, verificar que el planner elide particiones irrelevantes:
EXPLAIN SELECT * FROM events
WHERE created_at BETWEEN '2025-01-01' AND '2025-01-31';

-- Debe aparecer solo events_2025_01, no todas las particiones
-- Si aparecen todas: revisar que el filtro use la columna de partición exacta

-- Manejar partition por hash (distribución uniforme sin semántica temporal):
CREATE TABLE users_sharded (id BIGSERIAL, email TEXT)
PARTITION BY HASH (id);

CREATE TABLE users_sharded_0 PARTITION OF users_sharded
  FOR VALUES WITH (modulus 4, remainder 0);
CREATE TABLE users_sharded_1 PARTITION OF users_sharded
  FOR VALUES WITH (modulus 4, remainder 1);
-- ... etc para remainder 2 y 3
\end{lstlisting}

\section[Citus: Sharding sobre PostgreSQL]{\emj{🏗️}~Citus: Sharding sobre PostgreSQL}

Citus es una extensión de PostgreSQL (ahora parte de Microsoft) que convierte PG en una BD distribuida.
Un nodo coordinador recibe queries; las distribuye a workers que contienen fragmentos de datos.

\begin{teoriabox}{Conceptos clave de Citus}
\begin{itemize}
  \item \textbf{Distributed table:} particionada en shards distribuidos entre workers por una \emph{distribution column}
  \item \textbf{Reference table:} replicada en todos los workers (para tablas pequeñas en joins frecuentes: catálogos, configuración)
  \item \textbf{Local table:} solo en el coordinador (metadatos del sistema)
  \item \textbf{Colocation:} shards de dos tablas con la misma distribution column en el mismo worker → joins locales, sin red
\end{itemize}
\end{teoriabox}

Citus convierte PostgreSQL en una base de datos distribuida añadiendo un nodo \textbf{coordinador} que recibe todas las queries y las divide entre los \textbf{workers}.
La clave de diseño más importante es elegir la \textbf{distribution column}: la columna por la cual se fragmentan los datos entre workers.
Para un SaaS multi-tenant, \texttt{tenant\_id} es la elección natural: garantiza que todos los datos de un tenant estén en el mismo worker (colocation), haciendo que las queries por tenant sean locales y no requieran comunicación de red entre workers.

\begin{lstlisting}[caption={Setup básico de Citus}]
-- Instalar extensión:
CREATE EXTENSION citus;

-- Agregar workers al coordinador:
SELECT citus_add_node('worker1', 5432);
SELECT citus_add_node('worker2', 5432);

-- Distribuir tabla por tenant_id (SaaS típico):
SELECT create_distributed_table('orders', 'tenant_id');
SELECT create_distributed_table('items', 'tenant_id');  -- colocated con orders

-- Reference table (catálogo de productos: misma en todos los workers):
SELECT create_reference_table('products');

-- Queries single-tenant van directo al shard correcto (rápido):
SELECT * FROM orders WHERE tenant_id = 42;

-- Queries cross-tenant se ejecutan en paralelo en todos los workers:
SELECT tenant_id, COUNT(*) FROM orders GROUP BY tenant_id;
\end{lstlisting}

\section[Foreign Data Wrappers (FDW)]{\emj{🔗}~Foreign Data Wrappers (FDW)}

FDW (Foreign Data Wrappers) es el mecanismo de PostgreSQL para acceder a fuentes de datos externas como si fueran tablas locales.
PostgreSQL incluye \texttt{postgres\_fdw} para conectar con otras instancias de PG, y existen wrappers de terceros para MySQL, MongoDB, S3, archivos CSV, Redis, y más.
El uso más común en arquitecturas modernas es la \textbf{federación}: consultar datos de múltiples bases en una sola query SQL, o implementar migraciones incrementales donde la base nueva consulta datos de la vieja durante el período de transición.
El planner intenta hacer ``push-down'' de filtros y proyecciones al servidor remoto para minimizar los datos transferidos por red.

FDW permite acceder a fuentes externas (otra base PG, MySQL, S3, CSV) como si fueran tablas locales.
Útil para federación de datos y migraciones incrementales.

\begin{lstlisting}[caption={postgres\_fdw — acceso a otra base PG}]
-- En el servidor local:
CREATE EXTENSION postgres_fdw;

CREATE SERVER remote_orders
  FOREIGN DATA WRAPPER postgres_fdw
  OPTIONS (host 'remote-pg.example.com', port '5432', dbname 'orders_db');

CREATE USER MAPPING FOR current_user
  SERVER remote_orders
  OPTIONS (user 'readonly_user', password 'secret');

-- Importar todas las tablas de un schema remoto:
IMPORT FOREIGN SCHEMA public
  FROM SERVER remote_orders
  INTO local_remote_schema;

-- O crear tabla foránea individual:
CREATE FOREIGN TABLE remote_customers (
  id    BIGINT,
  email TEXT,
  name  TEXT
) SERVER remote_orders OPTIONS (schema_name 'public', table_name 'customers');

-- Usar como tabla local (push-down de filtros si es posible):
SELECT * FROM remote_customers WHERE id = 42;
\end{lstlisting}

\section[CockroachDB y Vitess: alternativas]{\emj{🚀}~CockroachDB y Vitess: alternativas}

\begin{teoriabox}{Comparativa de sistemas distribuidos}
\begin{itemize}
  \item \textbf{CockroachDB:} NewSQL compatible con PostgreSQL wire protocol. Distribuye datos automáticamente, transacciones distribuidas con Raft consensus. Operacional más simple que Citus pero overhead de latencia por consensus en cada write.
  \item \textbf{Vitess:} sharding para MySQL (usado por YouTube, Slack). Middleware: no modifica MySQL internamente. Maneja connection pooling, query routing, resharding online. Compatible con MySQL 8.
  \item \textbf{YugabyteDB:} open-source, compatible con PG y Cassandra API. DocDB storage engine + Raft. Bueno para geografía (multi-region).
  \item \textbf{PlanetScale:} Vitess as a service. Branching de schema (como git para la BD).
\end{itemize}
\end{teoriabox}

\begin{entrevistabox}
\begin{enumerate}
  \item ¿Cuándo tiene sentido sharding y cuándo es prematuro?
  \item ¿Cuál es la diferencia entre partitioning y sharding?
  \item En Citus, ¿qué es colocation y por qué importa para el rendimiento de joins?
  \item ¿Qué es una reference table en Citus? ¿Cuándo la usarías?
  \item Explica el problema de hot shard y cómo elegir una buena distribution column.
  \item ¿Cuándo usarías CockroachDB en lugar de PostgreSQL + Patroni?
\end{enumerate}
\end{entrevistabox}
